\chapter{Introduction}
Here we will responds: 
\begin{enumerate}
    \item What question I am trying to answer?\\
        Can the joint analysis of GW and EM signals encode relevant information about the existance of the phase transition inside a neutron star?
    \item Why is this question open an relevant?\\
        The posible confirmation or dismiss of quark-gluon plasma inside neutron stars could help us stablish the structure of neutron stars and test theoretical results in that regime.
    \item What will I do to contribute to this question?\\
        This thesis is just a prove of concept and a first approch to use the two different signals from a merger to try to identify posible encode information about the nature of the matter in the merger.
\end{enumerate}
\section[Dense matter and neutron-star equation of state]{Dense matter and neutron-star equation of state}
This section concentrates on the context of the second question. Here we will have a overview of the known equations of state, what they acount for and their limitations.
\section[Binary neutron-star mergers as multimessenger laboratories]{Binary neutron-star mergers as multimessenger laboratories}
This section is the state of the art of current efforts to use compact object mergers, in particular neutron star mergers, as test for theoretical develoment.
\section[Open problems: phase transitions and nuclear-physics uncertainties]{Open problems: phase transitions and nuclear-physics uncertainties}
This section as continuation of the first one states why is so difficult to calculate the equation of state of matter inside a neutron star.
\section[Objectives and research questions]{Objectives and research questions}
This section focus on the first question, it states what question am i answering.
Are combine GW and EM be usefull for testing theretical results?
\section[Structure of the thesis]{Structure of the thesis}
This is and overall review of what will i do in this thesis (answer to question 3)
